\documentclass[10pt]{article}

\usepackage[a4paper,top=13mm,bottom=12mm,left=18mm,right=18mm]{geometry}
\usepackage[T1]{fontenc}
\usepackage[utf8]{inputenc}
\usepackage[spanish]{babel}
\usepackage{lmodern}
\usepackage{microtype}
\usepackage[table]{xcolor}
\usepackage{tabularx}
\usepackage{array}
\usepackage{booktabs}
\usepackage{ragged2e}
\usepackage{eso-pic}

\definecolor{wine}{HTML}{8A2636}
\definecolor{ink}{HTML}{25282D}
\definecolor{muted}{HTML}{777A80}
\definecolor{rulegray}{HTML}{C9C9C9}
\definecolor{phasefill}{HTML}{F3F3F3}

\color{ink}
\setlength{\parindent}{0pt}
\setlength{\parskip}{2.15pt}
\setlength{\tabcolsep}{3.5pt}
\renewcommand{\arraystretch}{1.20}
\pagestyle{empty}

\newcommand{\proposalheading}[1]{%
  \vspace{1.8mm}%
  {\color{wine}\sffamily\fontsize{9.8}{11.2}\selectfont\bfseries
  \MakeUppercase{#1}}\par
  \vspace{0.55mm}%
}

\newcolumntype{F}{>{\raggedright\arraybackslash\color{wine}\sffamily\bfseries\fontsize{7.8}{9}\selectfont}p{25mm}}
\newcolumntype{A}{>{\raggedright\arraybackslash\fontsize{7.8}{9}\selectfont}X}
\newcolumntype{E}{>{\raggedright\arraybackslash\fontsize{7.8}{9}\selectfont}p{50mm}}

\begin{document}
\AddToShipoutPictureBG*{%
  \AtPageUpperLeft{\raisebox{-\height}{\color{wine}\rule{\paperwidth}{4mm}}}%
}

\begin{minipage}[t][0.98\textheight][t]{\textwidth}
\vspace{3.5mm}

\begin{center}
  {\fontsize{24}{27}\selectfont Planificación de Horarios Universitarios\par}
  \vspace{2.0mm}
  {\color{muted}\sffamily\fontsize{8.8}{10.2}\selectfont\bfseries
  ESTADO DEL ARTE Y REPRODUCIBILIDAD EN LA\par
  PLANIFICACIÓN DE HORARIOS BASADA EN PLANES DE ESTUDIO\par}
  \vspace{2.0mm}
  {\color{wine}\rule{36mm}{0.8pt}}
\end{center}

\proposalheading{Problema y objetivo}
{\fontsize{9.4}{11.2}\selectfont\justifying
La planificación de horarios universitarios asigna clases a periodos y aulas bajo restricciones de profesorado, espacios y planes de estudio. La calidad del horario también depende de la capacidad de las aulas, la distribución semanal, la compactación de las clases y la estabilidad del aula. Este trabajo se centra en el problema \textbf{Curriculum-Based Course Timetabling (CB-CTT)} y en la formulación pública UD2 de ITC-2007. Su objetivo es revisar los métodos actuales y comprobar si pueden reproducirse resultados recientes obtenidos con software de código abierto.\par}

\proposalheading{Alcance de la revisión}
{\fontsize{9.2}{11}\selectfont\justifying
La revisión abarcará investigaciones sobre CB-CTT publicadas entre 2015 y septiembre de 2026, además de los trabajos fundacionales del benchmark. Las búsquedas en Scopus, Web of Science e IEEE Xplore se ampliarán mediante rastreo de citas. Los estudios se compararán según su formulación, conjunto de datos, método, evidencia, entorno computacional y disponibilidad del código. La revisión incluirá métodos exactos y basados en restricciones, metaheurísticas, enfoques híbridos y técnicas asistidas por aprendizaje automático. ITC-2019 servirá como referencia más reciente porque incorpora la asignación de estudiantes a grupos y restricciones institucionales más completas.\par}

\vspace{0.7mm}
\noindent\begin{tabularx}{\linewidth}{@{}F A E@{}}
\toprule[0.65pt]
\rowcolor{phasefill}
\color{wine}\sffamily\bfseries FAMILIA &
\color{wine}\sffamily\bfseries COBERTURA DE LA REVISIÓN &
\color{wine}\sffamily\bfseries LÍNEA RECIENTE\\
\midrule[0.45pt]
Métodos exactos y de restricciones & MIP, CP, CP-SAT, descomposición y cotas & Construcción basada en modelos y descomposición\\
\midrule[0.35pt]
Metaheurísticas & SA, búsqueda tabú, VNS/VND, ALNS e hiperheurísticas & Simulated Annealing con LectureKick [5]\\
\midrule[0.35pt]
Métodos híbridos & Construcción asistida por solvers y mejora local & Inicialización con CP-SAT seguida de VND [6]\\
\midrule[0.35pt]
Métodos con aprendizaje & Ajuste basado en características, selección de algoritmos y control adaptativo & Espacio de instancias y selección de algoritmos [3]\\
\bottomrule[0.65pt]
\end{tabularx}

\proposalheading{Estudio de reproducibilidad}
{\fontsize{9.2}{11}\selectfont\justifying
El trabajo práctico reproducirá una selección de experimentos de Ceschia et al. [5]. Su repositorio público incluye el código fuente, las instancias del benchmark y la configuración de parámetros utilizada en el estudio. Compilaremos y ejecutaremos el solver sobre las \textbf{19 instancias no redundantes de ITC-2007 UD2}, compararemos los valores de la función objetivo y los tiempos de ejecución con los resultados publicados, y documentaremos el entorno y las discrepancias observadas. Se excluye \texttt{comp11} porque los métodos competitivos alcanzan de forma consistente un coste cero. La instancia \texttt{comp15} es equivalente a \texttt{comp03} bajo UD2 [1].\par}

\proposalheading{Contribución y alcance esperados}
{\fontsize{9.2}{11}\selectfont\justifying
El artículo combinará una revisión crítica de CB-CTT con la evaluación de la reproducibilidad de un solver reciente de código abierto. La discusión abordará la solidez de la evidencia disponible, la transparencia al comunicar los resultados y los límites de aplicar los resultados de ITC-2007 a la planificación universitaria actual.\par}

\proposalheading{Referencias}
{\fontsize{7.2}{8.3}\selectfont\raggedright
[1] Ceschia, Di Gaspero, and Schaerf, ``Educational timetabling: Problems, benchmarks, and state-of-the-art results,'' 2023. \quad
[2] Abdipoor et al., ``Metaheuristic approaches for university timetabling problems,'' 2023. \quad
[3] De la Rosa-Rivera et al., ``Algorithm selection for educational timetabling,'' 2021.\par
[4] Gu et al., ``From integer programming to machine learning: A survey on solving university timetabling problems,'' 2025. \quad
[5] Ceschia et al., ``Revisiting local search for curriculum-based course timetabling,'' 2026. \quad
[6] Nguyen, ``Adaptive penalty-guided VND with CP-SAT for curriculum-based course timetabling,'' 2026. \quad
[7] Di Gaspero, McCollum, and Schaerf, \textit{The Second International Timetabling Competition (ITC-2007): Curriculum-Based Course Timetabling}, 2007.}

\vfill
{\color{rulegray}\rule{\linewidth}{0.55pt}}\vspace{0.6mm}
{\color{muted}\sffamily\fontsize{7.1}{8.2}\selectfont
Metodología y Técnicas de Investigación \textbar{} Práctica 1 \textbar{} 2026-2027
\hfill Propuesta de investigación}

\end{minipage}
\end{document}
